\section{Introduction}\label{sec:intro}

In the \emph{shortest common superstring problem}, the input is a finite collection $\mathcal S$ of finite strings, and the task is to find a shortest string containing every member of $\mathcal S$ as a contiguous substring. Write $\OPT(\mathcal S)$ for its minimum possible length. The strings are represented explicitly, with encoded symbol labels and string boundaries. Let $N$ be the total bit length of this representation and $L$ the sum of the string lengths measured in symbols. Each symbol and boundary occupies at least one bit, so $L,|\mathcal S|\le N$; operations on symbol labels are charged for their encoded lengths.

\begin{theorem}\label{thm:main}
There is a deterministic algorithm, polynomial in the total encoded input size, that on every finite collection $\mathcal S$ of explicitly represented ordinary strings outputs a common superstring $T$ satisfying
\[
\len{T}\le 2\OPT(\mathcal S).
\]
\end{theorem}

\paragraph{History and significance.}
Tarhio and Ukkonen proposed the factor-$2$ conjecture for the maximum-overlap Greedy procedure~\cite{tarhio1988}. Blum, Jiang, Li, Tromp, and Yannakakis established a polynomial-time $3$-approximation~\cite{blum1994}; Sweedyk later obtained a $5/2$ guarantee~\cite{sweedyk1999}. Subsequent guarantees include $2+11/23$ due to Mucha~\cite{mucha2013} and $(14+\sqrt{67})/9<2.466$ due to Englert, Matsakis, and Vesel\'{y}~\cite{emv2023}. An August 2026 preprint of Chukhin, Kulikov, Mihajlin, and Smal reports a $7/3$-approximation~\cite{cycle2026}. \Cref{thm:main} answers affirmatively the question of whether a polynomial-time factor-$2$ approximation exists.

Many earlier approximation algorithms use the distance graph on the input strings, whose edge costs are the lengths of prefixes left after maximum suffix--prefix overlap. A minimum-weight cycle cover gives a computable lower bound, but its cycles must then be opened and joined while controlling the extra length~\cite{blum1994}. Breslauer, Jiang, and Jiang use rotations of periodic strings to control overlaps~\cite{breslauer1997}; Mucha sharpens this analysis through lexicographically extremal primitive rotations~\cite{mucha2013}. These provide conceptual precedents for the periodic comparisons below; our lower bound and the graph to which we apply them are different.

Our graph representation comes from another formulation. Golovnev, Kulikov, and Mihajlin introduced a hierarchical graph in the bounded-length setting~\cite{golovnev2014}, and the Collapsing framework developed its all-substrings form~\cite{collapsing2020}. In this graph, appending or deleting a letter gives an edge between substring vertices, and a superstring is represented by a closed walk visiting the required strings. The separation between a balanced graph and the operations that connect it also appears in the hierarchical cycle-cover discussion of that framework~\cite[revised version, Section~7.2]{collapsing2020}. We use this representation to turn forced occurrence counts into a balanced graph and then connect it within the same cost bound.

The factor-$2$ question has also motivated conjectures about specific procedures. The classical Greedy conjecture concerns repeatedly merging a pair of maximum overlap, while the Collapsing conjecture concerns a prescribed normalization of doubled hierarchical-graph walks~\cite{collapsing2020}. A September 2026 preprint of Shibata reports a counterexample to the Greedy conjecture, with worst-case ratio at least $9/4$~\cite{greedy2026}. Our theorem concerns the algorithm proved here; neither of those prescribed procedures is assumed to attain factor $2$.

\paragraph{Proof overview.}
The vertices of our hierarchical graph are all substrings of the inputs, including the empty word. An up edge appends a letter and costs one; a down edge deletes the first letter and costs zero. A balanced multiset of edges has equal incoming and outgoing degree at each vertex. If its support is connected and visits the empty word and every required string, an Euler tour spells a common superstring whose length is its up-edge cost.

We first construct integer counts $m(s)$ no larger than the number of occurrences of $s$ in any common superstring. Their one-letter sum $W$ is therefore a lower bound on the optimum. Counts for longer words constrain those of their prefixes and suffixes. A further rule handles periodic words. Relative to a fixed periodic text, a bracketing word consists of a matching interval with a mismatching letter at each end. When the assigned count of a long word exceeds the contribution of its bracketing words, a prefix shorter by $k$ periods must have count at least $k+1$ above its own bracketing contribution. The resulting counts determine a balanced, possibly disconnected graph of cost $W$.

Next decompose that graph into closed walks. The letters appended during a walk, repeated in both directions, form a periodic text; each visited vertex is a finite interval of that text. Walks with the same periodic text can be rearranged, preserving every edge, into ordered walks that each traverse one least period. We call these walks \emph{layers}. A layer of period $p$ has cost $p$ and receives a further budget $p$, so all additional budgets together equal $W$. In the remaining case of the connection procedure, the count rule forces a shorter word to occur on a layer with a different aligned text. A visited interval agreeing with both texts gives an actual graph vertex at which to attach a connection.

The final task is to arrange these connections without spending a budget twice. Some closed walks connect their source layers to a target using only the source budgets. Others require a contribution from the target, but these requests go only to a strictly larger period and are handled together when that period is processed. For each connection paid for only by its source layers, those layers can instead be joined internally while saving one period of budget. The remaining source-funded connections may form directed cycles. On each cycle, a maximal lexicographic rotation yields a contact word no longer than the target's period. Omitting that target's outgoing connection and joining its layers internally therefore pays for a round trip from the contact to the empty word. This connects the base graph at additional cost at most $W$ and completes the factor-$2$ bound.

\paragraph{Organization.}
\Cref{sec:counts} constructs the counts and their balanced graph. \Cref{sec:layers} develops the layer representation and the periodicity facts. \Cref{sec:connections} gives the connection operations, \Cref{sec:processing} processes groups of equal primitive text, and \Cref{sec:cycles} resolves the remaining cycles. \Cref{sec:algorithm} proves polynomial time and extracts the output string.
